\section{Methods}
\label{sec:methods}

\begin{figure}[t]
\centering
\resizebox{\linewidth}{!}{%
\begin{tikzpicture}[
    font=\normalsize,
    >=Latex,
    node distance=1.4cm and 1.7cm,
    box/.style={
        draw,
        rounded corners=3pt,
        line width=0.8pt,
        align=center,
        minimum height=1.05cm,
        minimum width=2.8cm,
        inner sep=6pt
    },
    input/.style={
        box,
        fill=green!8,
        minimum width=2.6cm
    },
    process/.style={
        box,
        fill=blue!7,
        minimum width=3.2cm,
        minimum height=1.35cm
    },
    channelK/.style={
        box,
        fill=blue!18,
        minimum width=2.8cm
    },
    channelR/.style={
        box,
        fill=green!13,
        minimum width=2.8cm
    },
    test/.style={
        box,
        fill=red!6,
        minimum width=4.1cm,
        minimum height=0.85cm
    },
    arrow/.style={
        ->,
        line width=0.85pt,
        draw=black!70
    }
]

% Title
\node[font=\large\bfseries] (title) at (5.3, 2.2)
    {Channel Decomposition Pipeline};

% Inputs
\node[input] (mem) at (0, 1.1) {Member\\texts};
\node[input] (nonmem) at (0, -1.1) {Non-member\\texts};

% Activations and decomposition
\node[process] (act) at (3.2, 0) {LLM layer $\ell$\\activations};
\node[process, font=\small] (bcd) at (6.9, 0) {Contrastive principal\\component analysis\\decomposition};

% Channels
\node[channelK] (sk) at (10.2, 1.1) {$S_K$\\\textbf{Knowledge}};
\node[channelR] (sr) at (10.2, -1.1) {$S_R$\\\textbf{Recall}};

% Tests
\node[test, font=\small] (geom) at (14.2, 1.7) {Geometry $\cos(d_K,d_R)$};
\node[test, font=\small] (erase) at (14.2, 0.55) {Erasure dissociation test};
\node[test, font=\small] (rank) at (14.2, -0.55) {Intrinsic rank from SVD};
\node[test, font=\small] (transfer) at (14.2, -1.7) {Cross-architecture transfer};

% Arrows
\draw[arrow] (mem.east) -- (act.west);
\draw[arrow] (nonmem.east) -- (act.west);
\draw[arrow] (act.east) -- (bcd.west);

\draw[arrow] (bcd.east) -- (sk.west);
\draw[arrow] (bcd.east) -- (sr.west);

\draw[arrow] (sk.east) -- (geom.west);
\draw[arrow] (sk.east) -- (erase.west);
\draw[arrow] (sr.east) -- (rank.west);
\draw[arrow] (sr.east) -- (transfer.west);

\end{tikzpicture}%
}
\caption{
Channel decomposition pipeline. Given member and non-member texts,
we extract layer-$\ell$ activations from a language model and apply contrastive
principal component analysis to obtain two subspaces, a knowledge channel $S_K$ and a recall
channel $S_R$ (each \emph{channel} is a contrastively estimated low-rank activation subspace, see \S\ref{sec:methods:bcd}). We then evaluate the decomposition with four diagnostics.
Geometry tests whether knowledge and recall directions are distinct. Erasure
tests whether the distinction has causal consequences. SVD estimates whether
the signal is concentrated or distributed. Cross-architecture transfer tests
whether the geometry is model-specific or shared.
}
\label{fig:behavioral-channel-decomposition}
\end{figure}

We probe mean-pooled residual stream activations at one layer of each fine-tuned model and decompose each activation in two complementary ways, along membership-knowledge and recall directions ($\dK$, $\dR$), and separately into an SAE reconstruction and reconstruction residual. The experiments test three claims. First, the standard SAE reconstruction preserves the membership-discriminating direction so that detection accuracy survives reconstruction. Second, the SAE feature dictionary explains the residual signal, so that targeted feature suppression removes membership recoverability from the residual. Third, the residual-over-reconstruction ordering generalizes across architectures within models that pass the validity gate of \S\ref{sec:methods:disclosures}. We separate two evidence sources. The controlled Pythia experiments use single-corpus member and non-member splits with known fine-tuning labels. Three configurations (Pythia-6.9B member-only, Pythia-6.9B mixed-data, and Pythia-12B mixed-data) are the primary tests of the first two claims. Models other than Pythia probe the third claim where the validity gate is satisfied and remain exploratory where it is not. We frame the residual experiments as a recoverability test, not a claim that the residual contains a single interpretable privacy feature or a complete causal localization of membership information.

\subsection{Fine-tuning setup with ground-truth membership labels}
\label{sec:methods:setup}

We fine-tune base models on a controlled subset of OpenWebText following the fine-tuning membership setup of \citet{meeus2025sok}. Membership labels are
known for this fine-tuning stage. Members are documents used for fine-tuning, and non-members are same-distribution OpenWebText documents from a disjoint
partition. Figure~\ref{fig:dark-sub-overview} in Appendix~\ref{app:models} shows the full SAE training and evaluation pipeline. This same-corpus construction reduces temporal-shift and vocabulary-imbalance confounds~\citep{duan2024membership,meeus2025sok}. Each fine-
tuned model uses 1{,}000 members and 1{,}000 non-members, all truncated to a common token length. We evaluate nine models from 1B to 12B parameters across
seven architecture
families~\citep{biderman2023pythia,black2021gpt,zhang2022opt,penedo2023refinedweb,jiang2023mistral,grattafiori2024llama3,yang2024qwen2}. Pythia-1B also
has epoch and pretraining checkpoints. Other fine-tuned models are analyzed after five epochs.

\subsection{Separating membership knowledge from recall behavior}
\label{sec:methods:bcd}

For the channel decomposition, we compute mean-pooled residual stream activations at each transformer layer. We use \emph{channel} as a shorthand for a
contrastively estimated low-rank activation subspace. The knowledge channel contrasts member and non-member documents, while the recall channel contrasts
high-recall and low-recall member documents.

We estimate the knowledge subspace $\SK$ with contrastive principal component analysis~\citep{abid2018exploring} on member and non-member activation
populations. We eigendecompose the contrastive covariance $C_{\text{mem}} - \alpha\, C_{\text{non}}$ with $\alpha = 1.0$ and keep the leading $n_c = 10$
eigenvectors as the subspace basis. We estimate the recall subspace $\SR$ analogously, contrasting high-recall and low-recall members. Recall labels use a
loss-based proxy that splits members at the median per-token cross-entropy loss. We validate this proxy against ROUGE-L extraction scores for each model
(Appendix~\ref{app:bcd_details}).

We summarize the relationship between $\SK$ and $\SR$ using cosine similarity between their leading directions $\dK$ and $\dR$, principal angles between
the subspaces, and cross-validated detector AUROC on each projection.

\subsection{Where membership signal survives SAE reconstruction}
\label{sec:methods:sae_recon}

We next test whether SAE reconstruction preserves membership evidence in the activation or leaves it recoverable from the reconstruction residual. We use the channel decomposition to define the membership direction, then test membership detection on the original activation, the SAE reconstruction, and the reconstruction residual.

\paragraph{Threat model and reconstruction split.}
We adopt a white-box internal-audit threat model in the sense of \citet{carlini2023quantifying}. The auditor can read layer activations and SAE
reconstructions, and has labeled member and same-distribution non-member documents for direction estimation and detector calibration. For each model, the
measured object is the mean-pooled residual stream activation $\mathbf{h}$ at the SAE target layer. We encode and decode $\mathbf{h}$ through the SAE to
obtain the reconstruction $\mathbf{h}' = \text{dec}(\text{enc}(\mathbf{h}))$ and the reconstruction residual $\mathbf{r} = \mathbf{h} - \mathbf{h}'$. We
then measure membership detection AUROC with logistic regression on $\mathbf{h}$, $\mathbf{h}'$, and $\mathbf{r}$.

If detection drops on $\mathbf{h}'$ but is preserved in $\mathbf{r}$, then the standard SAE reconstruction does not preserve the membership evidence
available in the original activation. As stated in the overview, the residual experiments are a recoverability test.

As a control for SAE training bias, we train a separate SAE on both member and non-member activations and repeat the experiment. This tests whether the
reconstruction drop is explained by a dictionary trained only on member activations. A fully unsupervised-direction variant of the pipeline is future
work.

\paragraph{Targeting the known membership direction.}
We compare two SAE training configurations. The \emph{member-only SAE} is trained only on member documents. The \emph{mixed-data SAE} is trained on the union of member and non-member documents and is therefore label-blind by construction. We test whether the residual membership signal is exhausted by the known direction $\dK$. We fine-tune a pre-trained mixed-data SAE with an additional
penalty on reconstruction error along $\dK$.
\begin{equation}
\mathcal{L} = \mathcal{L}_{\text{recon}} + \lambda_1 \|\mathbf{z}\|_1 + \lambda_{d_K} \frac{1}{N}\sum_{i=1}^{N} \bigl((\mathbf{h}_i - \mathbf{h}'_i)
\cdot \hat{\dK}\bigr)^2,
\label{eq:dk_penalty}
\end{equation}
\noindent where $\mathcal{L}_{\text{recon}}$ is the standard SAE reconstruction loss, $\lambda_1 \|\mathbf{z}\|_1$ is the baseline L1 sparsity penalty, and the third
term penalizes the squared projection of the residual onto the unit-normalized membership direction $\hat{\dK}$. Fine-tuning hyperparameters and the held-out audit and disjoint-evaluation partition split are in Appendix~\ref{app:fresh_probe_v2}.

\subsection{Whether SAE features carry membership and recall}
\label{sec:methods:ablation}

We next test whether membership and recall directions are encoded by identifiable SAE features. This matters for feature-level interventions. If the relevant channel is spanned by a small set of SAE decoder
directions, then ablating those features should affect the corresponding membership or recall outcome.

For each model, we first choose the SAE analysis layer. Before training any SAE, we run a layer-wise logistic membership detector on raw mean-pooled
residual stream activations under a fixed five-fold split, and select the layer with highest cross-validated AUROC. Per-model selected layers are listed
in Appendix Table~\ref{tab:model_details}.

At the selected layer, we call an SAE feature a classifier feature if its activation differs between members and non-members at false discovery
rate of $5\%$. We then ablate the top-$k$ classifier features in the SAE code, decode the modified code back into activation space, and measure both
membership detection and verbatim extractability. This tests whether the features most associated with membership labels are also sufficient to control
the downstream membership and recall outcomes.

To test whether the selected SAE features span the channel they are meant to explain, let $\mathcal{F}$ be a selected set of SAE features. Each feature has a decoder vector, and those decoder vectors span a subspace of activation space. Let
$\mathbf{P}_{\mathcal{F}}$ denote the orthogonal projection onto that feature subspace. For a target channel subspace $\mathcal{S}$, we define the feature
sufficiency criterion as
\begin{equation}
\mathrm{FSC}(\mathcal{S}, \mathcal{F})
=
\frac{\|\mathbf{P}_{\mathcal{F}}(\mathcal{S})\|_F^2}
 {\|\mathcal{S}\|_F^2},
\label{eq:fsc}
\end{equation}
\noindent where the numerator measures how much of the target channel is captured by the selected feature directions and the denominator is the channel's total size. Thus $\mathrm{FSC}$ is the fraction of the channel captured by the selected SAE features (1 means full coverage, 0 means no coverage). The full per-axis interpretation is in Appendix~\ref{app:fsc}.

\subsection{Interventions separate extraction from detection}
\label{sec:methods:dd}

We next use subspace erasure to test whether membership detection and verbatim extraction depend on separable activation directions. The intervention removes one channel from the
activation and measures what changes downstream. If erasing a channel suppresses verbatim extraction without suppressing membership detection, then
extraction and detection depend on different activation directions.

For a channel with orthonormal basis $\mathbf{V}$, erasure removes the component of the activation lying in that channel. Given a mean-pooled residual
stream activation $\mathbf{h}$, we define the channel-erased activation as
\begin{equation}
\mathbf{h}_{\setminus \mathbf{V}}
=
\mathbf{h} - \mathbf{V}\mathbf{V}^{\top}\mathbf{h},
\label{eq:subspace_erasure}
\end{equation}
\noindent where $\mathbf{V}\mathbf{V}^{\top}\mathbf{h}$ is the projection of $\mathbf{h}$ onto the channel subspace, and $\mathbf{h}_{\setminus \mathbf{V}}$ is the
remaining activation after that projection is removed.

We apply this operation separately to the knowledge basis $\mathbf{V}_K$ and the recall basis $\mathbf{V}_R$. After each intervention, we measure membership detection AUROC and verbatim extractability with ROUGE-L. Sensitivity to subspace dimension and the full per-condition statistical procedure are in Appendix~\ref{app:dd_full} and Appendix~\ref{app:dd_details}.

\subsection{Residual geometry repeats across architectures}
\label{sec:methods:transfer}

We next test whether the same channel geometry appears across model architectures. Three geometric questions follow. We measure shared channel geometry by principal angles between the knowledge subspaces of model pairs at matching hidden dimension. We measure how spread out member activations are with the participation ratio. We measure whether the member-versus-non-member shift lies inside the dominant variance directions with the cross-architecture concentration $C_k$ of the membership direction in the top-$k$ principal-component basis.

Each measure has a clear intuitive reading. Small principal angles indicate that two models place their knowledge subspaces in the same region. A high participation ratio indicates spread-out member variance, while a low ratio indicates compression into a few directions. $C_k$ near 1 means the membership shift lives along the dominant variance directions, while $C_k$ near 0 means it lies on lower-variance or off-subspace directions. Computational mechanics, the SVD identity, the formulae, the non-linear classifier comparison, and the architecture-family grouping rationale are in Appendix~\ref{app:cross_arch_readings}, Appendix~\ref{app:transfer}, Appendix~\ref{app:dimensionality_details}, and Appendix~\ref{app:nonlinear}.

\subsection{Validity gate for quantitative claims}
\label{sec:methods:disclosures}

We interpret cross-architecture comparisons only for SAE runs satisfying predefined validity criteria.

The validity gate decides which model-SAE settings we interpret quantitatively. We interpret an SAE decomposition only when reconstruction is reliable enough for the claim being made. The strict criterion requires reconstruction cosine $> 0.90$ and mean L0 below 60\% of the dictionary size. The controlled Pythia result satisfies this strict criterion, and a more permissive criterion at reconstruction cosine $> 0.85$ is used only for the qualitative cross-architecture ordering. Per-model reconstruction outcomes, exclusions, and the dictionary-expansion plus L1 settings used to clear the gate are in Appendix~\ref{app:qwen2_pilot} and Appendix~\ref{app:per_model_hps_detail}.

Standard L1-regularized ReLU SAE training is known to be per-initialization unstable near the sparsity phase transition~\citep{bricken2023monosemanticity, templeton2024scaling}. Across the cross-architecture sweep at the literature-default L1, some random initializations train into a sparse basin (a small fraction of features active per token) and others into a dense basin (most features active per token). The validity gate filters dense-basin training runs by construction, since their reconstruction cosine falls below the threshold. Per-seed evidence and the mechanism are in Appendix~\ref{app:sae_l1_instability}.

The dictionary-corpus choice matters because it controls whether the decomposition can see the evaluation labels. The mixed-data SAE is trained on the union of the member and non-member pools, so the dictionary is label-blind by construction. To rule out the related concern that SAE training documents overlap with evaluation documents, we retrain the Pythia-6.9B mixed-data SAE on an OpenWebText partition disjoint from the evaluation pool (\S\ref{sec:results:robustness}, Appendix~\ref{app:corpus_disjoint}). Member and non-member partitions are also IID by construction (same OpenWebText corpus, different documents). Recent work shows that pre-training MIAs without IID splits are largely confounded by distribution shift~\citep{maini2024dataset, duan2024membership}. The controlled fine-tuning regime here is in scope for those guidelines.

The detector scope matters because the primary effect size depends on it. The membership score is score$_K$, the projection of an activation onto the fitted knowledge direction $\dK$. AUROC is the main separability metric in the body, with low-FPR operating points in Appendix~\ref{app:tpr_paraphrase}. We also re-fit $\dK$ on a held-out partition to test partition sensitivity (Appendix~\ref{app:heldout_dk_protocol}). The paraphrase robustness check uses word-order shuffling, and cross-corpus replication is future work. Sparsity-range scope of the geometric and feature-alignment analyses is in Appendix~\ref{app:gemma_stretch}.